
\chapter*{Danh mục từ viết tắt}
\addcontentsline{toc}{chapter}{Danh mục từ viết tắt}
\markboth{Danh mục từ viết tắt}{Danh mục từ viết tắt}

\renewcommand{\arraystretch}{1.35}
\begin{longtable}{p{3.2cm} p{10.8cm}}
\hline
\textbf{Từ viết tắt} & \textbf{Giải thích} \\
\hline
\endfirsthead
\hline
\textbf{Từ viết tắt} & \textbf{Giải thích} \\
\hline
\endhead
\hline
\endfoot

% --- Nhóm Thuật toán & Giáo dục ---
AP      & Achievement Points — Điểm thành tích tích lũy của người dùng. \\
CEFR    & Common European Framework of Reference for Languages — Khung tham chiếu ngôn ngữ châu Âu thống nhất, phân loại trình độ từ A1 đến C2. \\
EF      & Easiness Factor — Hệ số dễ trong thuật toán SM-2, phản ánh mức độ quen thuộc với một từ vựng. \\
ELO     & Elo Rating System — Hệ thống xếp hạng kỹ năng cạnh tranh được áp dụng trong các trận đấu PvP. \\
HP      & Hit Points — Chỉ số sinh tồn của thú ảo trong các trận đối kháng. \\
SDT     & Self-Determination Theory — Lý thuyết tự quyết, nền tảng tâm lý học cho thiết kế động lực trong gamification. \\
SM-2    & SuperMemo 2 — Thuật toán lặp lại ngắt quãng xác định lịch ôn tập tối ưu dựa trên EF và Interval. \\
SRS     & Spaced Repetition System — Hệ thống lặp lại ngắt quãng, chiến lược học thuật tối ưu hóa việc ghi nhớ dài hạn. \\
TTL     & Time To Live — Thời gian tồn tại của một bản ghi trong bộ nhớ đệm trước khi bị xóa tự động. \\
XP      & Experience Points — Điểm kinh nghiệm, đơn vị đo lường tiến trình tăng trưởng của người dùng. \\

\hline
% --- Nhóm Kiến trúc & Lập trình ---
API     & Application Programming Interface — Giao diện lập trình ứng dụng, chuẩn hóa cách các thành phần phần mềm trao đổi dữ liệu. \\
CDN     & Content Delivery Network — Mạng phân phối nội dung, giúp giảm độ trễ tải tài nguyên tĩnh theo vị trí địa lý. \\
CRUD    & Create, Read, Update, Delete — Bốn thao tác cơ bản trên dữ liệu. \\
DI      & Dependency Injection — Kỹ thuật tiêm phụ thuộc, cơ chế quản lý vòng đời đối tượng trong ASP.NET Core. \\
DTO     & Data Transfer Object — Đối tượng truyền dữ liệu, tách biệt tầng trình bày khỏi tầng miền. \\
ERD     & Entity-Relationship Diagram — Lược đồ thực thể quan hệ, công cụ mô hình hóa cơ sở dữ liệu. \\
HTTP    & HyperText Transfer Protocol — Giao thức truyền siêu văn bản, nền tảng giao tiếp của World Wide Web. \\
HTTPS   & HyperText Transfer Protocol Secure — Phiên bản bảo mật của HTTP sử dụng mã hóa TLS/SSL. \\
JSON    & JavaScript Object Notation — Định dạng trao đổi dữ liệu nhẹ, dễ đọc và phổ biến trong REST API. \\
JWT     & JSON Web Token — Tiêu chuẩn mã thông báo truyền thông tin xác thực an toàn giữa các bên. \\
MVC     & Model-View-Controller — Mẫu kiến trúc phân tách giao diện, logic nghiệp vụ và dữ liệu. \\
ORM     & Object-Relational Mapping — Kỹ thuật ánh xạ đối tượng-quan hệ, cho phép tương tác với cơ sở dữ liệu thông qua ngôn ngữ lập trình hướng đối tượng. \\
REST    & Representational State Transfer — Phong cách kiến trúc thiết kế Web API dựa trên nguyên tắc phi trạng thái. \\
SDK     & Software Development Kit — Bộ công cụ phát triển phần mềm, thư viện và tài liệu do nhà cung cấp cung cấp. \\
SQL     & Structured Query Language — Ngôn ngữ truy vấn có cấu trúc dùng để quản lý cơ sở dữ liệu quan hệ. \\
UML     & Unified Modeling Language — Ngôn ngữ mô hình hóa thống nhất, bộ ký hiệu tiêu chuẩn biểu diễn thiết kế phần mềm. \\
URL     & Uniform Resource Locator — Địa chỉ định vị tài nguyên trên mạng. \\

\hline
% --- Nhóm Giao diện & Trải nghiệm ---
UI      & User Interface — Giao diện người dùng, tổng thể các thành phần trực quan người dùng tương tác. \\
UX      & User Experience — Trải nghiệm người dùng, chất lượng tổng thể của quá trình tương tác với sản phẩm. \\

\hline
% --- Nhóm Trò chơi ---
PvE     & Player versus Environment — Chế độ đối kháng giữa người chơi và hệ thống (Gym Leader). \\
PvP     & Player versus Player — Chế độ đối kháng trực tiếp giữa hai người chơi thời gian thực. \\
TTS     & Text-to-Speech — Công nghệ tổng hợp giọng nói từ văn bản, được sử dụng để tạo câu hỏi nghe hiểu. \\

\hline
% --- Nhóm Kiểm thử & Vận hành ---
P95     & 95th Percentile — Bách phân vị thứ 95 của phân phối độ trễ phản hồi, chỉ số đánh giá hiệu năng hệ thống dưới tải cao. \\
VU      & Virtual User — Người dùng ảo, đơn vị mô phỏng tải đồng thời trong kiểm thử hiệu năng k6. \\

\end{longtable}
